\documentclass[11pt]{article}

\usepackage[margin=1in]{geometry}
\usepackage{booktabs}
\usepackage{hyperref}

\title{Haplux: Reference-Neutral Auditing of Genomic Model Stability Across Haplotype Contexts}
\author{Ilgaz Mehmetoglu}
\date{\today}

\begin{document}

\maketitle

\begin{abstract}
Haplux is an early-stage research system for testing whether a genomic model's
prediction for the same focal variant is stable across different human haplotype
backgrounds. The central object of study is a paired effect,
\[
  \Delta_h = f(x_{h,\mathrm{ALT}}) - f(x_{h,\mathrm{REF}}),
\]
where both model inputs are constructed within the same local haplotype context
\(h\). Haplux aims to make this comparison reproducible by separating source
provenance, coordinate handling, haplotype reconstruction, model input policy,
and effect summarization. The first intended research contribution is not a new
predictive model, but an auditing framework for measuring when existing models
produce context-dependent variant effects.
\end{abstract}

\section{Motivation}

Genomic prediction models are commonly queried by replacing a reference allele
with an alternate allele inside a reference-derived sequence window. That
workflow is simple, but it can hide an assumption: the measured effect may depend
on the surrounding haplotype background used to construct the sequence. If nearby
variants, indels, phase, or ancestry-associated sequence backgrounds alter model
inputs, then a single reference-context effect can understate uncertainty or
misrepresent the range of model behavior.

Haplux investigates this problem by treating the focal variant effect as a
paired, context-indexed quantity. The same REF and ALT allele states are compared
within each compatible haplotype, and stability is summarized across those
paired effects.

\section{Research Question}

The primary research question is:

\begin{quote}
For a fixed focal variant and a fixed genomic prediction model, how stable is the
model's predicted REF-to-ALT effect across compatible human haplotype contexts?
\end{quote}

This decomposes into several operational questions:

\begin{itemize}
  \item Which haplotype contexts are compatible with a focal variant query?
  \item How should local REF and ALT model inputs be constructed when either
        allele may already be observed on a source haplotype?
  \item How much of the predicted effect variation is attributable to nearby
        sequence context rather than to the focal allele alone?
  \item Which model input policies, such as fixed-length centering or padding,
        change the measured stability?
  \item Can context-dependent model behavior be reported without making clinical
        interpretation claims?
\end{itemize}

\section{Current System}

The current Haplux prototype establishes the following technical base:

\begin{itemize}
  \item a 0-based half-open internal coordinate convention with VCF-style
        1-based input conversion;
  \item validation that declared reference alleles match the supplied sequence;
  \item construction of matched REF and ALT inputs when either allele is
        observed in the supplied haplotype;
  \item source provenance for linear references, phased VCF haplotypes,
        pangenome paths, de novo assemblies, and raw sequence inputs;
  \item GA4GH refget-style sequence content identifiers for constructed model
        inputs;
  \item a provider boundary that reconstructs local diploid haplotypes from
        indexed FASTA and phased VCF or BCF files;
  \item deterministic development scorers used to test paired-effect mechanics;
  \item strict, versioned JSON reports for headless and interactive workflows.
\end{itemize}

\section{Methodological Direction}

Haplux should improve variant-effect analysis by making the following quantities
explicit and machine-readable:

\begin{table}[h]
\centering
\begin{tabular}{@{}ll@{}}
\toprule
Quantity & Purpose \\
\midrule
Source provenance & Records where each context came from. \\
Sequence content identity & Detects when distinct sources yield identical inputs. \\
Observed allele state & Distinguishes observed and constructed REF/ALT inputs. \\
Model input policy & Records trimming, centering, padding, and fixed-length rules. \\
Paired effect & Compares ALT and REF inside the same haplotype context. \\
Stability summary & Reports spread, sign agreement, and skipped contexts. \\
\bottomrule
\end{tabular}
\end{table}

The near-term methodological improvement is to move from toy scorers to real
model adapters while preserving this audit trail. Each adapter should declare
its input policy, model identity, output target, software version, and failure
mode. Fixed-length models are especially important because trimming and padding
can change the local sequence seen by the model.

\section{Hypotheses}

Haplux can be used to test the following working hypotheses:

\begin{enumerate}
  \item Some focal variants have model effects whose sign or magnitude changes
        across valid haplotype contexts.
  \item Nearby phased variants and indels explain part of this effect spread.
  \item Fixed-length model input policies can introduce measurable differences
        in paired effects, especially near window boundaries and indels.
  \item Context-stability summaries provide useful uncertainty signals that are
        absent from single-reference-context variant scoring.
\end{enumerate}

\section{Evaluation Plan}

A practical evaluation should proceed in stages:

\begin{enumerate}
  \item Validate coordinate and allele handling on synthetic fixtures with known
        expected sequences.
  \item Reconstruct phased haplotypes at compact public-data loci and compare
        paired effects across chromosome copies.
  \item Add at least one real sequence model adapter with a documented
        fixed-length input policy.
  \item Compare reference-only effects against multi-haplotype effect
        distributions.
  \item Report effect range, standard deviation, sign agreement, skipped-context
        reasons, and model-input content identifiers.
\end{enumerate}

\section{Boundaries}

Haplux is research software. It should not provide clinical interpretation,
diagnosis, or patient-specific recommendations. Early claims should be limited
to reproducible model auditing: whether a model's predicted variant effect is
stable across compatible haplotype contexts under a stated model input policy.

\section{Next Improvements}

The next implementation milestones are:

\begin{enumerate}
  \item expand the model adapter contract for real inference backends;
  \item formalize fixed-length sequence policies in experiment request schemas;
  \item add normalization and left-alignment for VCF alleles;
  \item resolve overlapping or compound nearby variants without silent guessing;
  \item support cohort-scale batching across focal variants;
  \item add richer context-attribution summaries once model outputs are real.
\end{enumerate}

\end{document}
